\scanpage{47}
(on est ramené à voir que sur un corps $\unsure{\mathbb F_p|F_p|k_p}$ ayant \uline{au \deleted{\unsure{mo|moins}}moins trois éléments}, pour un vectoriel $M$ de rg $2$ sur ce corps, le \deleted{réunion} complémentaire de la réunion de deux droites engendre $M$ tout entier \ldots).
% slanted
\margin{à expliciter l'argument \ldots}

On a associé au module $M$ libre de rang $2$ sur $\Lambda=\Lambda_n=\mathbb Z/n\mathbb Z$, le torseur sous $Z'_{\Lambda_n}=Z'_n=(\mathbb Z/n\mathbb Z)^*/{\pm1}$\deleted{\unsure{,}}
\[
\pi=(\mathrm{det}\,M)^*/(\pm1_\Lambda),\tag*{(114)}
\]
à l'aide de celui-ci on réconstitue la ``droite mod $\pm1$'' \deleted{par} tordant
\[
L^{\prime\deleted{\unsure{*}}}=(\mathrm{det}\,M)/\{\pm1\}\simeq(\mathbb Z/n\mathbb Z)/\{\deleted{\unsure{p}}\pm1\}\mathbin{\wedge}_{Z'_n}\pi\ ;\tag*{(115)}
\]
on a une opération de $Z'_n$ sur $L'$, et $\pi$ en tant que $Z'_n$-ensemble s'identifie à une \deleted{\unsure{$Z$}} orbite de $Z'_n$ dans $L'$ -- savoir \inserted{à} l'unique orbite de cardinal \unsure{égale|égal} à $\varphi(n)/2=\mathrm{card}\ Z'_n$ :
\[
\pi\subset L'.\tag*{(116)}
\]
D'autre part, on associe à $M$ l'ens
\[
S=M^*/\{\pm\mathrm{id}_M\}\tag*{(117)}
\]
dont on veut saisir les structures essentielles, et les liens entre ceux-ci.
% archive sheet number: 46
